% Generated draft: requires mathematical and source review.
\documentclass[11pt]{article}
\usepackage[margin=25mm]{geometry}
\usepackage{fontspec}
\setmainfont{Latin Modern Roman}
\usepackage{amsmath,amssymb,mathtools,unicode-math}
\setmathfont{Latin Modern Math}
\usepackage{enumitem,needspace,xurl,hyperref}
\setlength{\parindent}{0pt}
\setlength{\parskip}{5pt}
\setlength{\emergencystretch}{3em}
\newcommand{\N}{\mathbb{N}}
\newcommand{\Z}{\mathbb{Z}}
\newcommand{\Q}{\mathbb{Q}}
\newcommand{\R}{\mathbb{R}}
\newcommand{\C}{\mathbb{C}}
\newcommand{\F}{\mathbb{F}}
\newcommand{\A}{\mathbb{A}}
\newcommand{\PP}{\mathbb{P}}
\newcommand{\E}{\mathbb{E}}
\newcommand{\eps}{\varepsilon}
\newcommand{\dd}{\,\mathrm{d}}
\begin{document}
\section*{UnsolvedMath ID: 30003366}
\textbf{Source problem code:} catalogue entry
\subsection{Definitions and mathematical statement}
A closed manifold is a compact manifold without boundary. For a Riemannian manifold \((X,g_X)\), strictly negative sectional curvature means
\[
K_{g_X}(\sigma)<0
\]
for every two-dimensional tangent plane \(\sigma\subseteq T_xX\) and every \(x\in X\).

The supplied question is: for every pair of closed negatively curved Riemannian manifolds \(M,N\) satisfying
\[
\pi_1(M)\cong\pi_1(N),
\]
must their \(k\)-frame bundles be isomorphic, for every admissible \(k\)? The source does not specify the meaning of frame-bundle isomorphism or the precise range of \(k\), so this question does not yet determine a unique mathematical assertion.

For clarity, one conventional meaning of the \(k\)-frame bundle of an \(n\)-dimensional Riemannian manifold \(X\), for an integer \(1\leq k\leq n\), is the orthonormal tangent-frame bundle
\[
V_k(TX,g_X)
=
\left\{(x,v_1,\ldots,v_k):
 x\in X,\ v_i\in T_xX,\
 g_X(v_i,v_j)=\delta_{ij}\text{ for }1\leq i,j\leq k
\right\},
\]
with projection \(p_X(x,v_1,\ldots,v_k)=x\). Its fiber is the real Stiefel manifold \(V_k(\R^n)\). Another convention uses all ordered linearly independent tangent vectors instead of orthonormal ones. Neither convention is explicitly selected in the record.

Under the orthonormal convention and an additional common-dimension hypothesis \(\dim M=\dim N=n\), a bundle-isomorphism interpretation would ask whether there exist maps
\[
\Phi:V_k(TM,g_M)\longrightarrow V_k(TN,g_N),
\qquad f:M\longrightarrow N,
\qquad p_N\circ\Phi=f\circ p_M,
\]
that constitute an isomorphism in a specified bundle category. A topological bundle isomorphism requires appropriate homeomorphisms; a smooth bundle isomorphism requires appropriate diffeomorphisms. Any further structure to be preserved must also be specified. An isomorphism of total spaces alone is a different possible target. These are conditional interpretations, not hypotheses or conclusions supplied by the source. The source also does not require the maps to induce a prescribed isomorphism of fundamental groups.

The parenthetical assertion that the answer is known for \(k=1,2\) is retained as a supplied claim, for values within the applicable frame range. Its precise category and hypotheses require verification alongside the intended general question.
\subsection{Short English statement}
The question asks whether isomorphic fundamental groups force closed negatively curved manifolds to have isomorphic \(k\)-frame bundles. The source reports affirmative results for \(k=1,2\) and labels the general question open. The frame convention, admissible dimensions, and meaning of bundle isomorphism remain unspecified, so a definitive formal statement requires human review.
\subsection{Sources}
\begin{itemize}
\item Supplied local UnsolvedMath registry record: permanent ID \(30003366\), title ``Frame Bundle Rigidity of Negatively Curved Manifolds,'' source problem code \texttt{catalogue entry}, in the Oberwolfach Reports Open Problems collection. No individual proposer is identified; the proposed year is recorded as 2017.
\item Supplied report citation: \emph{Mini-Workshop: Spaces and Moduli Spaces of Riemannian Metrics (2018)}, Oberwolfach Reports, pp.~133--166, DOI \texttt{10.4171/owr/2017/3}. Supplied link: \url{https://doi.org/10.4171/owr/2017/3}. The citation's parenthetical year 2018 differs from the report identifier and recorded proposed year 2017; this discrepancy is unresolved.
\item Supplied literature source: ``Oberwolfach Report 3/2017,'' \url{https://ems.press/doi/pdf/10.4171/OWR/2017/3}.
\item Supplied literature source: Gromov, \emph{Topological rigidity problems}, \url{https://www.ihes.fr/~gromov/wp-content/uploads/2018/08/177.pdf}.
\item Literature assessment dated exactly \texttt{2026-08-21}: the supplied curation labels the problem open, attributes the cases \(k=1,2\) to Gromov--Cheeger results, and states that no general theorem for higher frame bundles has been established. This is supplied dated curation derived from the OWR material, not an independently verified literature assessment. No precise theorem locator or separate Cheeger reference is supplied.
\item The registry labels the statement \texttt{corrected\_verified}, explaining that notation, typography, and sentence structure were repaired without changing mathematical content. This supplied editorial label does not establish mathematical correctness or independent source verification.
\end{itemize}
\end{document}
